\originalchapter{最小生成树与最小割}\label{ch:mst}
\sourceof{18.433 L9 的随机缩边; 生成树交换与贪心算法为编者补充}
\originalsection{割性质与圈性质}
设连通无向图每边有实权 $w(e)$. 最小生成树是总权最小的生成树. 图有限, 树的数量有限且至少一棵, 所以最小树存在; 边权可为负, 不需要最短路算法的非负假设.
\begin{theorem}[安全边]\label{thm:safeedge}
设边集 $A$ 已包含于某棵最小生成树. 若一个割不被 $A$ 中任何边穿过, 则该割中的一条最轻边 $e$ 可加入 $A$, 新集合仍包含于某棵最小生成树.
\end{theorem}
\begin{proof}
取包含 $A$ 的最小树 $T$. 若 $e\in T$ 已完成. 否则加 $e$ 后产生唯一圈, 圈沿 $e$ 跨割后必须有另一条边 $f$ 跨回. 因割不穿过 $A$, $f\notin A$. 又 $w(e)\le w(f)$, 用 $e$ 替换 $f$ 得一棵包含 $A\cup\{e\}$ 的树, 权不增加. 由原树最优性, 新树也是最小树.
\end{proof}
\begin{proposition}[圈排除]
若边 $e$ 在某个圈上严格比其他边重, 它不在任何最小生成树中. 若各边权互异, 最小生成树唯一.
\end{proposition}
\begin{proof}
树若含 $e$, 删它分两块, 原圈上必有另一边 $f$ 跨两块. 替换后权严格减少, 矛盾. 唯一性方面, 若两棵最小树 $T,T'$ 不同, 取其对称差中权最小的边 $e$, 可设它在 $T$ 中. 删除 $e$ 得割, $T'$ 中跨割的路含某条 $f\notin T$, 否则原 $T-e$ 已跨割. 最小性与权互异使 $w(e)<w(f)$, 在 $T'$ 中替换得更轻树, 矛盾.
\end{proof}
严格性不能省略. 等权三角形的任意两条边都是最小树, 所以圈上一条并列最重边仍可能属于最小树.
\originalsection{Kruskal 与 Prim}
Kruskal 先把边按权非递减排序, 初始 $A=\varnothing$, 若下一边端点处在当前森林不同分支就加入, 否则跳过. 它不产生圈. 加入时任选其一端所在分支形成割, 所有更轻的跨割边已被扫描; 若未加入, 当时两端已在同分支, 以后也仍在同分支, 不可能跨当前割. 因而当前边为该割的最轻边, 割不穿过 $A$, 安全边定理每步保持可扩为最小树. 最后森林极大, 连通输入使它为生成树, 所以最优. 并查集维护分支, 排序与合并可在 $O(m\log m)$ 时间完成.

Prim 从一点起, 维护已连成的顶点集 $S$ 及其树边 $A$, 每次加入 $S$ 到补集最轻边及其新端点. 割 $\partial S$ 不穿过 $A$, 所以每次安全. 连通性保证到全部顶点前割非空, 最后得到最小生成树. 邻接表与二叉堆实现为 $O((n+m)\log n)$, 矩阵实现为 $O(n^2)$.
\begin{example}
图有边 $ab:1,bc:2,cd:3,ac:4,bd:5,ad:6$. 求最小生成树, 并给出最优性证据.
\end{example}
\begin{solution}
按 Kruskal 顺序加入 $ab,bc,cd$, 已连接四点, 总权6. 非树边 $ac$ 的树路最大权2, $bd$ 的为3, $ad$ 的为3, 都不超过非树边自身权. 这给出最优性证书: 任何其他树的某条非树边加入当前树都不能通过换掉该圈上一边减重. 该证书的充分性将在下段通过固定的基本割逐次交换证明, 避免误用变化中的树路.
\end{solution}
上面的证书充分性应选以原候选树的路径为基准来作交换. 可直接用割证书避免中间树路径变化: 对候选树任一边 $f$, 删它所得基本割上每条非树边 $e$ 的原树路径包含 $f$, 故 $w(f)\le w(e)$. 因而每条树边都是其基本割的一条最轻边. 要把任意其他树改成候选树, 每次添加一个缺失候选边 $f$, 沿其他树中的相应圈删一条跨该基本割且不在候选树中的边 $e$, 权不增加. 有限次变为候选树, 因而原另一树权不小于候选树. 这一论证说明证书如何从局部比较提升到全局最优.
\originalsection{图的交换结构}
\begin{proposition}\label{prop:forestexchange}
若 $A,B$ 都为森林边集且 $|A|<|B|$, 则有 $e\in B\setminus A$ 使 $A\cup\{e\}$ 仍为森林.
\end{proposition}
\begin{proof}
若 $B$ 的每条边都连接 $A$ 的同一分支内部, 对每个 $A$ 分支的顶点集 $V_i$, 森林 $B[V_i]$ 的边数至多 $|V_i|-1$. 求和得到 $|B|\le\sum_i(|V_i|-1)=|A|$, 矛盾. 所以 $B$ 有边跨 $A$ 的两个分支, 加它不成圈.
\end{proof}
这是图拟阵的交换公理. 本册只在图的森林上证明它, 不把一般拟阵优化或多面体理论视为已覆盖. 贪心能正确, 正因为可以把局部选择与其他可行树通过交换联系起来.
\originalsection{随机缩边求最小割}
考虑连通无向多重图, 删除自环, 保留平行边的各自身份. Karger 算法反复在当前全部边中均匀随机选一条, 收缩其端点, 丢弃新自环, 直到只剩两个超级顶点. 剩下平行边形成原图一个割, 算法输出其大小. 收缩可能误删最小割, 因而单次只给概率保证.
\begin{theorem}\label{thm:karger}
在 $n\ge2$ 顶点连通图中, 一次随机缩边保留并输出某个固定最小割的概率至少为 $2/(n(n-1))$.
\end{theorem}
\begin{proof}
固定最小割大小为 $k\ge1$. 条件于先前没有收缩该割的边, 它在当前图仍存在且大小 $k$. 当前图的任何割对应原图的一个割, 所以最小割不小于 $k$, 每个超级顶点度至少 $k$. 若还剩 $r$ 点, 当前边数至少 $rk/2$, 下一步选到固定割边的概率至多 $k/(rk/2)=2/r$. 生存概率至少
\[
\prod_{r=n}^{3}\left(1-\frac2r\right)=\frac2{n(n-1)}.
\]
生存到两点时, 两超级顶点恰为所固定割的两側, 输出该最小割. $n=2$ 时空乘积为1, 与公式一致.
\end{proof}
独立重复 $L$ 次并取最小值, 失败概率至多 $(1-p)^L\le e^{-pL}$, $p=2/(n(n-1))$. 取 $L\ge n(n-1)\log(1/\eta)/2$ 可使失败概率不超过 $\eta$. 每次输出都是真实割, 不会输出低于真正最小值的虚假值. 非连通图的最小割为0, 应先以搜索识别, 不必进行缩边. 平行边不能合并后等概率选择, 否则证明的 $k/m$ 概率已被改变.
\originalsection{习题与解答}
\begin{exercise}\label{ex:mst1}
是否可以用 BFS 树代替最小生成树? 给出一个反例, 即使所有边权非负.
\end{exercise}
\begin{exercise}\label{ex:mst2}
证明任何最小生成树也是瓶颈最小生成树: 它最大边权在全部生成树中最小.
\end{exercise}
习题\ref{ex:mst1}的解答.
\begin{solution}
三角形 $ab:10,ac:10,bc:1$, 以 $a$ 为 BFS 根, 树边为 $ab,ac$, 总权20; 最小树为 $ab,bc$ 或 $ac,bc$, 总权11. BFS 最小化的是根到各点的边数距离, 与生成树总权不同.
\end{solution}
习题\ref{ex:mst2}的解答.
\begin{solution}
取最小树中的最重边 $e$, 删它得到基本割. 若有另一树所有边均严格轻于 $e$, 它必含一条跨割边 $f$, 用 $f$ 换 $e$ 就降低最小树总权, 矛盾. 所以其他树的最大边权至少为 $w(e)$.
\end{solution}
